\mwnewpart
\renewcommand\runningtitle{Section 1. Overview}
\chapter{Overview}
\sectioncontents

\section{What RetrOS/LT Is Made Of}\ixs{RetrOS/LT, parts of}

\begin{plainwords}
PC DOS 2.10, the DOS of the IBM Personal Computer in 1983, was made of
three kinds of program. Microsoft wrote the heart of it, the kernel and
the command processor, and most of the commands. IBM wrote the part that
works the computer's own hardware, and a set of commands of its own:
\typed{FDISK}, \typed{MODE}, \typed{BACKUP} and others. In 2018 Microsoft
published its part as source code, under a license that lets anyone build
on it. IBM's part has never been published. RetrOS/LT builds Microsoft's
part from that source, and writes IBM's part again, to behave as IBM's
did.
\end{plainwords}

\begin{tabularx}{\textwidth}{@{}P{1.6in}P{1.1in}Y@{}}
\toprule
\textbf{Part} & \textbf{From} & \textbf{In this repository}\\
\midrule
Kernel (\file{RETDOS.COM}) & Microsoft & \file{upstream/msdos}, with the
patches, and \file{src/IO.ASM}\\
\file{COMMAND.COM} & Microsoft & \file{upstream/msdos}, with the patches\\
SYSINIT & Microsoft & linked into \file{RETBIO.COM}\\
\typed{CHKDSK}, \typed{DEBUG}, \typed{EDLIN}, \typed{EXE2BIN}, \typed{FC},
\typed{FIND}, \typed{MORE}, \typed{PRINT}, \typed{RECOVER}, \typed{SORT},
\typed{SYS} & Microsoft & \file{upstream/msdos}, some with patches\\
\typed{FORMAT} & both & Microsoft's \file{FORMAT.ASM} and
\file{FORMES.ASM} with \file{src/IBMFOR.ASM}\\
\typed{LINK} & Microsoft & the program as released\\
\file{RETBIO.COM} & new & \file{src/BIOS.ASM}\\
Startup records & new & \file{boot/BOOT.ASM}, \file{boot/MBR.ASM}\\
\typed{ASSIGN}, \typed{BACKUP}, \typed{COMP}, \typed{DISKCOMP},
\typed{DISKCOPY}, \typed{FDISK}, \typed{GRAPHICS}, \typed{MODE},
\typed{RESTORE}, \typed{TREE} & new & \file{cmd/}\\
\file{ANSI.SYS} & new & \file{dev/ANSI.ASM}\\
Install & new, after RetrOS/x86's & \file{install/INSTALL.ASM}\\
\bottomrule
\end{tabularx}

\section{The Source It Starts From}\ixs{MS-DOS 2 source}

Microsoft's repository holds MS-DOS 1.25 and 2.0, and later 4.0; RetrOS/LT
pins it as the submodule \file{upstream/msdos}. Its \file{v2.0/source}
directory is not exactly 2.0: it is the kit Microsoft sent to computer
makers, and several files are of version 2.11, the release of March 1984
that joined PC DOS 2.10's corrections to MS-DOS 2.01. The kit has two
versions of some files (\file{DOSMAC.ASM} and \file{DOSMAC\_V211.ASM},
for one), and its \file{bin} directory holds programs built from it as
2.00. RetrOS/LT builds the 2.11 versions throughout, and reports 2.10.

The kit is incomplete in three ways that matter:

\begin{itemize}
\item \textbf{There is no BIOS.} Each maker wrote its own; the kit has a
  skeleton for one other computer, \file{SKELIO.ASM}. IBM's was never
  published. \file{src/BIOS.ASM} is RetrOS/LT's (Section~3).
\item \textbf{The kernel's console functions are missing.}
  \file{STDIO.ASM} includes a file \file{IO.ASM} that is not in the
  release. \file{src/IO.ASM} is a reconstruction (Section~4).
\item \textbf{IBM's commands are not there}, nor the part of
  \typed{FORMAT} for the IBM Personal Computer. They are written new
  (Section~5).
\end{itemize}

\section{The Patch Series}\ixs{patches}

RetrOS/LT's changes to Microsoft's files are kept as patches, applied in
the order of \file{patches/series}:

\srclist{series}{The patch series.}

\begin{tabularx}{\textwidth}{@{}P{0.5in}Y@{}}
\toprule
\patch{0001} & Builds with the 2.11 files: \file{DOSMAC} and \file{DOSSYM}
replaced by their \file{\_V211} versions; names that MASM 5 reserves
(\code{WAIT}, \code{DOSSEG}) renamed; a typo in \file{MSCODE.ASM}
corrected; SYSINIT set to IBM's way of starting \file{COMMAND.COM}; and
messages SYSINIT needs but the kit's \file{SYSIMES.ASM} lacked.\\
\patch{0002} & The RetrOS/LT identity: the banner, \typed{VER}, and DOS
version 2.10.\\
\patch{0003} & \typed{PRINT} from its 2.11 source.\\
\patch{0004} & \typed{FORMAT} for the IBM Personal Computer: IBM's switches,
and RetrOS/LT's system file names.\\
\patch{0005} & \typed{SYS} for the IBM Personal Computer: the names, and
the RetrOS/LT startup record.\\
\bottomrule
\end{tabularx}

Appendix~B lists the files each patch changes.

\section{The Version}\ixs{version, DOS}

PC DOS 2.10 reports major version 2, minor version 10, through interrupt
21h function 30h, and every program tests that. RetrOS/LT does the same:
\file{DOSSYM.ASM}'s \code{DOS\_MINOR\_VERSION} is 10 (\patch{0002}).
Microsoft's commands of the 2.11 kit refuse to run under a DOS older than
2.11 where they check at all; \typed{FORMAT}'s check is set to 2.10 (\patch{0004}).

\section{The Names}\ixs{RETBIO.COM}\ixs{RETDOS.COM}

PC DOS calls its system files \file{IBMBIO.COM} and \file{IBMDOS.COM};
MS-DOS calls them \file{IO.SYS} and \file{MSDOS.SYS}. RetrOS/LT's are
\file{RETBIO.COM} and \file{RETDOS.COM}, in IBM's pattern. Four places
know them: the startup record, \file{RETBIO.COM}'s loader for
\file{RETDOS.COM}, \typed{FORMAT} and \typed{SYS}.
